\begin{longtable}{p{.065\textwidth}p{.22\textwidth}p{.47\textwidth}p{.105\textwidth}}
\toprule 节点 & 主题 & 首次引入的技能与概念 & 作业与考试 \\\midrule\endhead
\multicolumn{4}{l}{\textbf{一、一阶微分方程}}\\
R1 & 自然增长与分离变量 & 指数增长及捕获建模; 增长率; 分离变量; 通解与特解; 合并积分常数; $\ln|y|$ 及消去; 补回损失解; 用常数满足初值 & -- \\
L1 & 方向场与存在唯一性 & 方向场、积分曲线、等斜线、漏斗; 隐式解; 导数无穷导致不能延续 & 发布作业1 \\
R2 & 方向场与分界曲线 & 分离轨线及解的极值 & -- \\
L2 & 数值方法 & Euler 方法 & -- \\
L3 & 线性方程与建模 & 一阶线性方程; 系统与信号视角; 银行账户; RC 电路; 常输入的分离变量解 & -- \\
R3 & Euler 方法及线性模型 & 混盐问题 & -- \\
L4 & 积分因子 & 齐次方程与零输入; 积分因子; 暂态; 扩散模型及耦合常数 & -- \\
R4 & 一阶线性方程 & 正弦输入 & -- \\
L5 & 复数与单位根 & 复数; 单位根 & 提交1; 发布2 \\
L6 & 复指数与正弦函数 & 复指数; 正弦的振幅、角频率及相位滞后 & -- \\
L7 & 指数及正弦输入的线性响应 & 一阶线性系统的指数/正弦响应; 正弦输入的复方程; 作业中的半衰期 & -- \\
R5 & 复数与复指数 & 原栏为空 & -- \\
L8 & 自治方程、相线与稳定性 & 自治方程; 相线; 稳定性; $e^{k(t-t_0)}$ 与 $ce^{kt}$ & 提交2; 发布3 \\
L9 & 线性与非线性 & 解不能延续的情形 & -- \\
R6 & 期中I复习 & 原栏为空 & 无独立题面 \\
-- & 期中I & 原栏为空 & 一小时考试 \\
\multicolumn{4}{l}{\textbf{二、二阶线性方程}}\\
R7 & 二阶方程的解 & 简谐振子; 初值; 齐次叠加 & -- \\
L11 & 模态与特征多项式 & 弹簧、质量及阻尼器; 一般二阶线性方程; 特征多项式; 实根解 & -- \\
L12 & 振动与阻尼 & 复根; 欠/过/临界阻尼; 复数替换与实解提取; 暂态; 根图 & -- \\
R8 & 二阶常系数齐次方程 & 一般正弦响应; 归一化解 & -- \\
L13 & 指数响应与弹簧驱动 & 受迫系统; 叠加; 指数响应公式; 复数替换; 正弦输入的正弦响应 & -- \\
R9 & 指数与正弦输入 & 原栏为空 & -- \\
L14 & 复增益与阻尼器驱动 & 增益、相位滞后及复增益 & 提交3; 发布4 \\
L15 & 算子、待定系数与共振 & 微分算子; 共振; 待定系数 & -- \\
R10 & 增益、共振与待定系数 & 原栏为空 & -- \\
L16 & 频率响应 & 频率响应 & -- \\
R11 & 频率响应 & 一阶频率响应 & -- \\
L17 & LTI 系统与 RLC 电路 & RLC 电路; 时间平移不变性 & 提交4; 发布5 \\
L18 & 工程应用 & 阻尼比 & -- \\
R12 & 期中II复习 & 原栏为空 & 无独立题面 \\
L19 & 期中II & 原栏为空 & 一小时考试 \\
\multicolumn{4}{l}{\textbf{三、Fourier 级数及变换响应}}\\
R13 & Fourier 级数入门 & 周期函数 & -- \\
L20 & Fourier 级数 & Fourier 级数; 正交性; Fourier 积分 & -- \\
L21 & Fourier 级数的运算 & 方波; 分段连续; 三角恒等式、线性组合及平移 & -- \\
R14 & Fourier 级数 & 不同周期 & -- \\
L22 & 周期解与共振 & 级数求导与积分; 谐波响应; Fourier 级数的振幅--相位表达 & -- \\
R15 & Fourier 谐波响应 & 原栏为空 & -- \\
L23 & 阶跃及脉冲 & 阶跃; $\delta$; 正则函数与奇异函数; 广义函数及导数 & 提交5; 发布6 \\
L24 & 阶跃响应与脉冲响应 & 单位阶跃/脉冲响应; 静止初值; 一阶、二阶响应 & -- \\
R16 & 阶跃、脉冲及其响应 & 原栏为空 & -- \\
L25 & 卷积 & 脉冲后的初值; 时间不变与 $D$、时间平移的交换; 卷积乘积; 初值响应 $w*q$ & -- \\
R17 & 卷积 & $\delta$ 为卷积单位元 & -- \\
L26 & Laplace 基本性质 & Laplace 变换; 收敛域; $\Lap[t^n]$; $s$ 平移; 正弦、余弦变换; 时域与变换域 & 提交6; 发布7 \\
L27 & Laplace 解方程 & $\Lap[\delta]$; 时间导数规则; 逆变换; 部分分式与遮盖法; 一阶非静止初值 & -- \\
R18 & Laplace 变换 & 用变换求单位阶跃响应 & -- \\
L28 & 二阶方程与配方 & $s$ 求导规则; 二阶方程 & -- \\
R19 & Laplace 方法II & 原栏为空 & -- \\
L29 & 极点图 & 权函数、传递函数及其变换关系; 时间平移; 极点及长期行为 & 提交7; 发布8 \\
L30 & 传递函数与频率响应 & 稳定性; 传递函数与增益 & -- \\
R20 & 期中III复习 & 原栏为空 & 无独立题面 \\
-- & 期中III & 原栏为空 & 一小时考试 \\
\multicolumn{4}{l}{\textbf{四、一阶方程组}}\\
L32 & 线性系统与矩阵 & 一阶线性系统; 消元; 矩阵; 反消元及伴随矩阵 & -- \\
R21 & 一阶线性系统 & 原栏为空 & -- \\
L33 & 特征值与特征向量 & 行列式; 特征值; 特征向量; 初值 & -- \\
R22 & 特征值与特征向量 & 解与轨道的区别 & -- \\
L34 & 复特征值或重根 & 特征值与系数; 复特征值; 重根; 缺陷情形与完整特征基 & 提交8; 发布9 \\
L35 & 线性系统相平面 & 迹--行列式平面; 稳定性 & -- \\
R23 & 线性相图 & 线性相图随参数变形 & -- \\
L36 & 正常模态与矩阵指数 & 矩阵指数; 解耦系统; 指数律 & -- \\
R24 & 矩阵指数 & 非齐次系统, 常输入 & -- \\
L37 & 非线性系统 & 非线性自治系统; 向量场; 相图; 平衡点; 平衡附近线性化; Jacobian 矩阵 & 提交9 \\
L38 & 平衡附近线性化及非线性单摆 & 非线性单摆; 飞机长周期振荡; Tacoma Narrows 桥 & -- \\
R25 & 自治系统 & 捕食--被捕食系统 & -- \\
L39 & 线性方法的局限 & 结构稳定; 极限环; 奇异吸引子 & -- \\
R26 & 课程复习 & 原栏为空 & 无独立题面 \\
-- & 综合期末 & 原栏为空 & 三小时考试 \\
\bottomrule
\end{longtable}
